% Generated by report/make_restated.py from theory/B2_effective_curvature.tex: do not edit (edit the note and rerun).
\begin{restated}{Lemma}{lem:B2-span}{B2}{1}{invariance of the feature span}
Let $S=\operatorname{span}\{x_1,\dots,x_n\}\subseteq\R^d$, let $S^\perp$ be its orthogonal complement, and write
$\theta=\theta_S+\theta_\perp$ for the orthogonal decomposition of a vector $\theta$ (similarly $v_S,v_\perp$).
For $\ell_2$-regularized logistic regression:
\begin{enumerate}[label=(\alph*)]
\item $\etastar\in S$, and $\fip(\theta)-\lambda\theta\in S$ for every $\theta$ and every $i$; in particular
$\fip(\theta)_\perp=\lambda\theta_\perp$ for all $i$, and $\Fp(\theta)_\perp=\lambda\theta_\perp$.
\item (Invariance) If $\theta_0\in S$, then, for gradient descent and for SAGA with the initialization
$z_i^{(0)}=\fip(\theta_0)$ of Proposition~5.9 (or with $z_i^{(0)}=0$), all iterates $\theta_t$, all memories $z_i^{(t)}$,
all directions $\omega_t$ and all error vectors $g_i^{(t)}$ lie in $S$, for every $t$ and every realization of the
indices. The same holds for the quadratic model of Lemma~\ref{lem:B2-saga-quadratic} whenever $H_iS\subseteq S$ for
every $i$ and $v_0,e_1^{(0)},\dots,e_n^{(0)}\in S$.
\item (The complementary component is autonomous) For an arbitrary $\theta_0$, the $S^\perp$-component of the
error $v_t=\theta_t-\etastar$ evolves as if $F$ were the quadratic $\frac\lambda2\norm{\theta_\perp}^2$. For gradient
descent, $v_{t,\perp}=(1-\gamma\lambda)^tv_{0,\perp}$. For SAGA, let $s_j(t)$ be the time at which the current memory was
computed, $z_j^{(t)}=f'_j(\theta_{s_j(t)})$ (so $s_j(0)=0$, $s_{i(t)}(t)=t-1$ and $s_j(t)=s_j(t-1)$ for $j\neq i(t)$); then
\[
v_{t,\perp}=(1-\gamma\lambda)\,v_{t-1,\perp}-\gamma\lambda\Big(\frac1n\sum_{j=1}^nv_{s_j(t-1),\perp}-v_{s_{i(t)}(t-1),\perp}\Big),
\]
and the error vectors satisfy
\[
g^{(t+1)}_{i,\perp}=\lambda\Big(v_{s_i(t),\perp}-\frac1n\sum_{j=1}^nv_{s_j(t),\perp}\Big),
\]
so that $\E[v_{t,\perp}]=(1-\gamma\lambda)^tv_{0,\perp}$, and the noise injected by SAGA in a unit direction $q\in S^\perp$ is
\[
q^\top C_{t+1}\,q\;=\;\lambda^2\,\Var_j\big(q^\top v_{s_j(t)}\big)
\;:=\;\lambda^2\cdot\frac1n\sum_{i=1}^n\Big(q^\top v_{s_i(t)}-\frac1n\sum_{j=1}^nq^\top v_{s_j(t)}\Big)^2 ,
\]
the empirical variance, over the memory, of the $q$-component of the stored errors.
\end{enumerate}
\end{restated}
